% Tutorial 9 solutions -- fix errors directly in this file.
\subsection*{Solutions --- Tutorial No. 9}
\medskip\noindent\textbf{Q.1  Eigenvalues, eigenvectors, multiplicities}\par\smallskip
AM = algebraic multiplicity (multiplicity as a root of $\det(A-\lambda I)$), GM = geometric multiplicity ($=\dim$ of the eigenspace).

\noindent\textbf{(1)} $A=\begin{pmatrix}1 & 4 \\ 2 & 3\end{pmatrix}$
\[ \det(A-\lambda I)=\left(\lambda - 5\right) \left(\lambda + 1\right)=0 \]
\[ \lambda=-1:\ \text{AM}=1,\ \text{GM}=1,\ \text{eigenvectors: }t\begin{pmatrix}-2 \\ 1\end{pmatrix}\ (t\neq0) \]
\[ \lambda=5:\ \text{AM}=1,\ \text{GM}=1,\ \text{eigenvectors: }t\begin{pmatrix}1 \\ 1\end{pmatrix}\ (t\neq0) \]

\noindent\textbf{(2)} $A=\begin{pmatrix}1 & -2 \\ -5 & 4\end{pmatrix}$
\[ \det(A-\lambda I)=\left(\lambda - 6\right) \left(\lambda + 1\right)=0 \]
\[ \lambda=-1:\ \text{AM}=1,\ \text{GM}=1,\ \text{eigenvectors: }t\begin{pmatrix}1 \\ 1\end{pmatrix}\ (t\neq0) \]
\[ \lambda=6:\ \text{AM}=1,\ \text{GM}=1,\ \text{eigenvectors: }t\begin{pmatrix}- \frac{2}{5} \\ 1\end{pmatrix}\ (t\neq0) \]

\noindent\textbf{(3)} $A=\begin{pmatrix}5 & 4 \\ 1 & 2\end{pmatrix}$
\[ \det(A-\lambda I)=\left(\lambda - 6\right) \left(\lambda - 1\right)=0 \]
\[ \lambda=1:\ \text{AM}=1,\ \text{GM}=1,\ \text{eigenvectors: }t\begin{pmatrix}-1 \\ 1\end{pmatrix}\ (t\neq0) \]
\[ \lambda=6:\ \text{AM}=1,\ \text{GM}=1,\ \text{eigenvectors: }t\begin{pmatrix}4 \\ 1\end{pmatrix}\ (t\neq0) \]

\noindent\textbf{(4)} $A=\begin{pmatrix}3 & 3 \\ 3 & 3\end{pmatrix}$
\[ \det(A-\lambda I)=\lambda \left(\lambda - 6\right)=0 \]
\[ \lambda=0:\ \text{AM}=1,\ \text{GM}=1,\ \text{eigenvectors: }t\begin{pmatrix}-1 \\ 1\end{pmatrix}\ (t\neq0) \]
\[ \lambda=6:\ \text{AM}=1,\ \text{GM}=1,\ \text{eigenvectors: }t\begin{pmatrix}1 \\ 1\end{pmatrix}\ (t\neq0) \]

\noindent\textbf{(5)} $A=\begin{pmatrix}14 & -10 \\ 5 & -1\end{pmatrix}$
\[ \det(A-\lambda I)=\left(\lambda - 9\right) \left(\lambda - 4\right)=0 \]
\[ \lambda=4:\ \text{AM}=1,\ \text{GM}=1,\ \text{eigenvectors: }t\begin{pmatrix}1 \\ 1\end{pmatrix}\ (t\neq0) \]
\[ \lambda=9:\ \text{AM}=1,\ \text{GM}=1,\ \text{eigenvectors: }t\begin{pmatrix}2 \\ 1\end{pmatrix}\ (t\neq0) \]

\noindent\textbf{(6)} $A=\begin{pmatrix}7 & 1 \\ -9 & 1\end{pmatrix}$
\[ \det(A-\lambda I)=\left(\lambda - 4\right)^{2}=0 \]
\[ \lambda=4:\ \text{AM}=2,\ \text{GM}=1,\ \text{eigenvectors: }t\begin{pmatrix}- \frac{1}{3} \\ 1\end{pmatrix}\ (t\neq0) \]
Since GM $<$ AM for some eigenvalue, $A$ is \emph{defective} (not diagonalizable).


\noindent\textbf{(7)} $A=\begin{pmatrix}6 & 2 \\ 2 & 3\end{pmatrix}$
\[ \det(A-\lambda I)=\left(\lambda - 7\right) \left(\lambda - 2\right)=0 \]
\[ \lambda=2:\ \text{AM}=1,\ \text{GM}=1,\ \text{eigenvectors: }t\begin{pmatrix}- \frac{1}{2} \\ 1\end{pmatrix}\ (t\neq0) \]
\[ \lambda=7:\ \text{AM}=1,\ \text{GM}=1,\ \text{eigenvectors: }t\begin{pmatrix}2 \\ 1\end{pmatrix}\ (t\neq0) \]

\noindent\textbf{(8)} $A=\begin{pmatrix}1 & 2 \\ 0 & 4\end{pmatrix}$
\[ \det(A-\lambda I)=\left(\lambda - 4\right) \left(\lambda - 1\right)=0 \]
\[ \lambda=1:\ \text{AM}=1,\ \text{GM}=1,\ \text{eigenvectors: }t\begin{pmatrix}1 \\ 0\end{pmatrix}\ (t\neq0) \]
\[ \lambda=4:\ \text{AM}=1,\ \text{GM}=1,\ \text{eigenvectors: }t\begin{pmatrix}\frac{2}{3} \\ 1\end{pmatrix}\ (t\neq0) \]

\noindent\textbf{(9)} $A=\begin{pmatrix}-2 & 5 & 4 \\ 5 & 7 & 5 \\ 4 & 5 & -2\end{pmatrix}$
\[ \det(A-\lambda I)=\left(\lambda - 12\right) \left(\lambda + 3\right) \left(\lambda + 6\right)=0 \]
\[ \lambda=-6:\ \text{AM}=1,\ \text{GM}=1,\ \text{eigenvectors: }t\begin{pmatrix}-1 \\ 0 \\ 1\end{pmatrix}\ (t\neq0) \]
\[ \lambda=-3:\ \text{AM}=1,\ \text{GM}=1,\ \text{eigenvectors: }t\begin{pmatrix}1 \\ -1 \\ 1\end{pmatrix}\ (t\neq0) \]
\[ \lambda=12:\ \text{AM}=1,\ \text{GM}=1,\ \text{eigenvectors: }t\begin{pmatrix}1 \\ 2 \\ 1\end{pmatrix}\ (t\neq0) \]

\noindent\textbf{(10)} $A=\begin{pmatrix}5 & 0 & 1 \\ 0 & -2 & 0 \\ 1 & 0 & 5\end{pmatrix}$
\[ \det(A-\lambda I)=\left(\lambda - 6\right) \left(\lambda - 4\right) \left(\lambda + 2\right)=0 \]
\[ \lambda=-2:\ \text{AM}=1,\ \text{GM}=1,\ \text{eigenvectors: }t\begin{pmatrix}0 \\ 1 \\ 0\end{pmatrix}\ (t\neq0) \]
\[ \lambda=4:\ \text{AM}=1,\ \text{GM}=1,\ \text{eigenvectors: }t\begin{pmatrix}-1 \\ 0 \\ 1\end{pmatrix}\ (t\neq0) \]
\[ \lambda=6:\ \text{AM}=1,\ \text{GM}=1,\ \text{eigenvectors: }t\begin{pmatrix}1 \\ 0 \\ 1\end{pmatrix}\ (t\neq0) \]
\medskip\noindent\textbf{Q.2  Eigenvalues of related matrices}\par\smallskip
If $A\mathbf x=\lambda\mathbf x$ then $A^{-1}\mathbf x=\lambda^{-1}\mathbf x$ (if $A$ is invertible), $A^{21}\mathbf x=\lambda^{21}\mathbf x$, $(A+I)\mathbf x=(\lambda+1)\mathbf x$, $(A^2+A+I)\mathbf x=(\lambda^2+\lambda+1)\mathbf x$, and $A^T$ has the same eigenvalues as $A$ (same characteristic polynomial).

\noindent\textbf{(1)} $A=\begin{pmatrix}1 & 1 & 0 \\ 0 & -1 & 0 \\ 0 & 0 & -1\end{pmatrix}$, $\det(A-\lambda I)=\left(\lambda - 1\right) \left(\lambda + 1\right)^{2}$.
\begin{center}\begin{tabular}{l|ccc}
 matrix & $\lambda_1$ & $\lambda_2$ & $\lambda_3$ \\ \hline
$A$ & $-1$ & $-1$ & $1$ \\
$A^{-1}$ & $-1$ & $-1$ & $1$ \\
$A^{21}$ & $-1$ & $-1$ & $1$ \\
$A+I$ & $0$ & $0$ & $2$ \\
$A^T$ & $-1$ & $-1$ & $1$ \\
$A^2+A+I$ & $1$ & $1$ & $3$ \\
\end{tabular}\end{center}

\noindent\textbf{(2)} $A=\begin{pmatrix}4 & 1 & 0 \\ 1 & 4 & 0 \\ 0 & 0 & 2\end{pmatrix}$, $\det(A-\lambda I)=\left(\lambda - 5\right) \left(\lambda - 3\right) \left(\lambda - 2\right)$.
\begin{center}\begin{tabular}{l|ccc}
 matrix & $\lambda_1$ & $\lambda_2$ & $\lambda_3$ \\ \hline
$A$ & $2$ & $3$ & $5$ \\
$A^{-1}$ & $\frac{1}{2}$ & $\frac{1}{3}$ & $\frac{1}{5}$ \\
$A^{21}$ & $2097152$ & $10460353203$ & $476837158203125$ \\
$A+I$ & $3$ & $4$ & $6$ \\
$A^T$ & $2$ & $3$ & $5$ \\
$A^2+A+I$ & $7$ & $13$ & $31$ \\
\end{tabular}\end{center}

\noindent\textbf{(3)} $A=\begin{pmatrix}2 & 1 & 1 \\ 1 & 2 & 1 \\ 1 & 1 & 2\end{pmatrix}$, $\det(A-\lambda I)=\left(\lambda - 4\right) \left(\lambda - 1\right)^{2}$.
\begin{center}\begin{tabular}{l|ccc}
 matrix & $\lambda_1$ & $\lambda_2$ & $\lambda_3$ \\ \hline
$A$ & $1$ & $1$ & $4$ \\
$A^{-1}$ & $1$ & $1$ & $\frac{1}{4}$ \\
$A^{21}$ & $1$ & $1$ & $4398046511104$ \\
$A+I$ & $2$ & $2$ & $5$ \\
$A^T$ & $1$ & $1$ & $4$ \\
$A^2+A+I$ & $3$ & $3$ & $21$ \\
\end{tabular}\end{center}
